\subsection{Conclusions}

The present work provides a combined numerical and experimental investigation of CB formation during shock-tube pyrolysis of 5\% CH$_4$ in argon at a post-reflected-shock pressure of P$_5$=4.5$\pm$0.5~atm and temperatures spanning 1800~K$<$T$_5$$<$2500~K. Multi-wavelength laser diagnostics were employed to measure the time-resolved mole fractions of CH$_4$, C$_2$H$_4$, and C$_2$H$_2$. The CB volume fraction, $f_v$, was determined using light extinction at 633~nm and 1064~nm, while collected TEM images were analyzed to extract the primary particle diameter, $d_p$, as a quantitative measure of CB morphology. Simulations using the CRECK mechanism accurately reproduced methane conversion and acetylene yields. The gas-phase chemistry solver was coupled with a sectional population balance model that accounts for CB inception, surface growth, and the reduction in H/C due to carbonization.

Simultaneous calibration of the model using both $f_v$ and $d_p$ provided strong constraints on inception, PAH adsorption, HACA surface growth, and carbonization rates, highlighting the importance of morphology data in avoiding non-unique calibrations based solely on $f_v$. The model reproduced the time history of $f_v$ and the final $d_p$ in good agreement with the measurements. Four stages of CB evolution were identified: (I) an induction period dominated by precursor formation, (II) rapid increase in $f_v$ driven by rising surface growth rates, (III) transition toward final yield, composition, and $d_p$ as surface growth wanes, and (IV) a coagulation-dominated stage in which primary particle diameters are chemically frozen while the agglomerate mobility diameter continues to increase.

The analysis showed that PAH adsorption dominates CB mass growth at lower temperatures and early times, whereas HACA becomes the primary contributor with increasing T$_5$. The model reproduced the bell-shaped temperature dependence of $f_v$ but overpredicted induction times at T$_5$$>$2200~K, suggesting the presence of missing high-temperature PAH formation pathways that would accelerate inception and PAH adsorption. TEM measurements revealed a decreasing trend in primary particle diameter with increasing T$_5$, while the model predicted a bell-shaped dependence. The discrepancy at 1984~K was attributed to possible late-time condensation of residual aromatics.

The maturity index, $\beta$, inferred from the extinction ratio at 633~nm and 1064~nm, exhibited strong similarities to the predicted H/C ratio of CB particles, which is used in the model as a measure of surface reactivity. Both $\beta$ and H/C decrease rapidly at early times before approaching near-asymptotic values, with the rate of the initial decline increasing with T$_5$.